\expandafter\ifx\csname CAPETINCLUDEMODE\endcsname\relax
\documentclass[11pt,a4paper]{article}
% preamble.tex — packages et commandes communes a toutes les fiches CAPET SII IE
\input{glyphtounicode}
\pdfgentounicode=1
\usepackage{cmap}
\usepackage[T1]{fontenc}
\usepackage[utf8]{inputenc}
\usepackage[french]{babel}
% babel[french] active a \begin{document} des raccourcis de ponctuation (:;!?)
% qui cassent la syntaxe to[...] de circuitikz (erreur "+ or - expected") ;
% on les desactive, ce qui ne change que l'espacement fin avant ces signes.
\AtBeginDocument{\shorthandoff{:;!?}}
\usepackage[a4paper, top=1.6cm, bottom=1.6cm, left=1.4cm, right=1.4cm]{geometry}
\setlength{\headheight}{14.5pt}
\usepackage{amsmath, amssymb, amsfonts}
\usepackage{siunitx}
\sisetup{locale=FR, output-decimal-marker={,}, inter-unit-product=\ensuremath{{}\cdot{}},
         group-minimum-digits=4}
% Unites non fournies par siunitx, declarees ici une fois pour toutes (fiches + exercices).
% Les gardes \ifdefined evitent tout conflit avec une redeclaration locale restee dans un
% fichier : ne pas retirer les gardes, ne pas redeclarer ces unites dans les fichiers.
\DeclareSIUnit{\jour}{j}
\ifdefined\var\else\DeclareSIUnit{\var}{var}\fi   % puissance reactive
\ifdefined\vad\else\DeclareSIUnit{\vad}{VAD}\fi   % puissance deformante
\ifdefined\tr\else\DeclareSIUnit{\tr}{tr}\fi      % tour (tr/min)
\usepackage{array, booktabs, multirow, tabularx}
\usepackage{xcolor}
\usepackage[most, breakable]{tcolorbox}
\usepackage{multicol}
\usepackage{enumitem}
\usepackage{mathtools}
\usepackage{lmodern}
% Refonte visuelle 2026 (direction "moderne coloré") : Lato pour le texte,
% maths sans-serif assorties via newtxsf (charge apres amsmath/amssymb).
\usepackage[default]{lato}
\usepackage{newtxsf}
\usepackage{microtype}
\usepackage{fancyhdr}
\usepackage{tikz}
\usetikzlibrary{positioning,arrows.meta}
\usepackage{circuitikz}
\usepackage{lastpage}

% ── Couleurs par domaine (palette modernisée 2026) ────────────────────
\definecolor{colElec}{RGB}{29, 78, 216}        % Électrotechnique — bleu vif
\definecolor{colElecLight}{RGB}{219, 234, 254} % fond bleu clair
\definecolor{colAuto}{RGB}{21, 128, 61}        % Automatique — vert
\definecolor{colMeca}{RGB}{194, 65, 12}        % Mécanique — orange
\definecolor{colTherm}{RGB}{185, 28, 28}       % Thermique — rouge
\definecolor{colQT}{RGB}{232, 244, 233}        % (hérité, plus utilisé par questiontype)
\definecolor{colJury}{RGB}{255, 248, 225}      % fond jury (crème)
\definecolor{colJuryFrame}{RGB}{202, 138, 4}   % bordure/bandeau jury (ambre)
\definecolor{colGrayLight}{RGB}{245, 245, 247} % fond gris neutre
\definecolor{colSeparator}{RGB}{180, 180, 180} % gris séparateur
\definecolor{colExo}{RGB}{8, 145, 178}          % Banque d'exercices — cyan (distinct des fiches)
% ── Encre et gris de la refonte ───────────────────────────────────────
\definecolor{colInk}{RGB}{23, 32, 42}          % texte courant (quasi-noir)
\definecolor{colSlate}{RGB}{83, 98, 118}       % gris ardoise (méta, en-têtes)
\definecolor{colRule}{RGB}{203, 213, 225}      % filets clairs
\definecolor{colJuryBg}{RGB}{255, 248, 225}    % fond crème jury
\definecolor{colPresence}{RGB}{83, 98, 118}    % bloc 5 (présence dans les sujets)
\AtBeginDocument{\color{colInk}}

% ── Style encadré exercice + corrigé (banque d'exercices) ─────────────
\tcbset{
  exoBox/.style={
    breakable,
    enhanced,
    colback=colExo!8!white,
    colframe=colExo,
    colbacktitle=colExo,
    coltitle=white,
    fonttitle=\bfseries\small,
    left=4pt, right=4pt, top=3pt, bottom=3pt,
    before skip=5pt, after skip=5pt,
  }
}

% ── Bascule enonce / corrige (banque d'exercices) ─────────────────────
% \CAPETCORRIGE est defini en ligne de commande (pdflatex "\def\CAPETCORRIGE{1}\input{...}")
% pour compiler la variante corriges d'un meme fichier source ; par defaut (non defini),
% la compilation produit la variante enonces. \CAPETSHOW{enonce}{corrige} choisit la branche.
%
% TROIS MODES (refonte 2026-09-20, manuel de revision) :
%   - aucun drapeau   -> imprime l'enonce            (variante enonces, INCHANGE)
%   - \CAPETCORRIGE   -> imprime le corrige          (variante corriges, INCHANGE)
%   - \CAPETMANUEL    -> imprime l'enonce ET met le corrige de cote pour l'annexe
%
% Le mode manuel stocke chaque corrige dans une macro (\csgdef), sans ecriture
% fichier ni re-tokenisation : le contenu (circuitikz, accents, %) traverse intact.
% L'annexe les rejoue ensuite dans l'ordre via \CAPETANNEXE.
\usepackage{etoolbox}

\newcounter{CAPETexo}          % numero continu d'exercice dans le manuel
\setcounter{CAPETexo}{0}
\newcounter{CAPETannexe}       % compteur de relecture en annexe

% Titre courant de la notion en cours (pose par \CAPETNOTION), sert d'intertitre annexe
\newcommand{\CAPETnotioncode}{}
\newcommand{\CAPETnotiontitre}{}

% \CAPETNOTION{code}{titre} : ouvre une notion dans le manuel (memorise le libelle,
% et marque le rang du premier exercice de la notion pour l'intertitre d'annexe).
\newcommand{\CAPETNOTION}[2]{%
  \renewcommand{\CAPETnotioncode}{#1}%
  \renewcommand{\CAPETnotiontitre}{#2}%
  \csgdef{CAPETnotionstart\the\numexpr\value{CAPETexo}+1\relax}{#1 -- #2}%
}

% ── \degree robuste ───────────────────────────────────────────────────
% Certains fichiers d'exercices utilisent \degree. Ni le preambule ni siunitx
% ne le definissent, et \providecommand local ne suffit pas : sous hyperref,
% \degree existe deja comme commande factice (\ERROR), donc \providecommand
% s'abstient et l'appel echoue a la compilation du manuel (erreur invisible
% en compilation isolee). On le (re)definit donc inconditionnellement ici,
% apres le chargement des packages, pour tout le corpus.
\AtBeginDocument{\providecommand{\degree}{}\renewcommand{\degree}{\ensuremath{^\circ}}}

% Libelle decoratif : les fichiers d'exercices utilisent aussi \CAPETSHOW dans
% leur \fancyhead et leur bloc de titre, pour afficher "(enonces)" ou "(corriges)".
% Ces appels-la ne sont PAS des exercices : en mode manuel ils ne doivent ni
% consommer un numero, ni alimenter l'annexe. On les reconnait a leur premier
% argument, strictement egal au libelle ci-dessous.
\newcommand{\CAPETLIBELLEENONCES}{(\'enonc\'es)}

\ifdefined\CAPETMANUEL
  % --- Mode manuel : enonce en place, corrige differe vers l'annexe ---
  \newcommand{\CAPETSHOW}[2]{%
    \def\CAPETargun{#1}%
    \ifx\CAPETargun\CAPETLIBELLEENONCES
      #1%            appel decoratif : rendu tel quel, hors numerotation
    \else
      \CAPETSHOWEXO{#1}{#2}%
    \fi
  }
  \newcommand{\CAPETSHOWEXO}[2]{%
    \stepcounter{CAPETexo}%
    \csgdef{CAPETcorrige\theCAPETexo}{#2}%
    \edef\CAPETtmpnum{\theCAPETexo}%
    #1%
    \par\vspace{2pt}%
    {\footnotesize\itshape\color{colExo}%
     \hyperref[capet:cor:\CAPETtmpnum]{Corrige de l'exercice \CAPETtmpnum\ $\rightarrow$ p.~\pageref{capet:cor:\CAPETtmpnum}}%
     \label{capet:eno:\CAPETtmpnum}}%
  }
  % \CAPETANNEXE : rejoue tous les corriges captures, dans l'ordre, avec intertitres
  \newcommand{\CAPETANNEXE}{%
    \setcounter{CAPETannexe}{0}%
    \loop\ifnum\value{CAPETannexe}<\value{CAPETexo}%
      \stepcounter{CAPETannexe}%
      \edef\CAPETcurnum{\theCAPETannexe}%
      \ifcsname CAPETnotionstart\CAPETcurnum\endcsname
        \par\medskip
        {\normalsize\bfseries\color{colExo}\csuse{CAPETnotionstart\CAPETcurnum}}%
        \par\nobreak\smallskip
      \fi
      \begin{tcolorbox}[exoBox, breakable,
        title={Corrige de l'exercice \CAPETcurnum\
               \hfill\normalfont\footnotesize
               \hyperref[capet:eno:\CAPETcurnum]{enonce p.~\pageref{capet:eno:\CAPETcurnum}}}]%
      \label{capet:cor:\CAPETcurnum}%
      \csuse{CAPETcorrige\CAPETcurnum}%
      \end{tcolorbox}%
    \repeat
  }
\else
  % --- Modes historiques, comportement strictement inchange ---
  \newcommand{\CAPETSHOW}[2]{\ifdefined\CAPETCORRIGE #2\else #1\fi}
  \newcommand{\CAPETANNEXE}{}
\fi

% ── Difficulté 1 à 3 (étoiles pleines, amssymb) ───────────────────────
\newcommand{\difficulte}[1]{%
  \ifcase#1\relax
  \or $\bigstar$\;niveau 1/3%
  \or $\bigstar\bigstar$\;niveau 2/3%
  \or $\bigstar\bigstar\bigstar$\;niveau 3/3%
  \fi
}

% ── Style encadré extrait de document technique (banque d'exercices) ──
\tcbset{
  docTech/.style={
    breakable,
    enhanced,
    colback=colGrayLight,
    colframe=colSeparator,
    fonttitle=\bfseries\small\itshape,
    coltitle=black,
    colbacktitle=colGrayLight,
    left=4pt, right=4pt, top=3pt, bottom=3pt,
    before skip=4pt, after skip=4pt,
  }
}

% ── Style tcolorbox de section principale ────────────────────────────
% Refonte 2026 : bandeau titre coloré arrondi + fond très pâle, sans cadre.
\tcbset{
  ficheSection/.style={
    breakable, enhanced jigsaw,
    frame hidden, boxrule=0pt, arc=1.8mm,
    colback=#1!4!white,
    colbacktitle=#1,
    coltitle=white,
    fonttitle=\sffamily\bfseries,
    left=7pt, right=7pt, top=5pt, bottom=5pt,
    toptitle=3pt, bottomtitle=3pt,
    before skip=10pt, after skip=8pt,
  }
}

% ── Style encadré question type ──────────────────────────────────────
% Carte blanche, bandeau sombre, liseré très clair.
\tcbset{
  questiontype/.style={
    breakable, enhanced jigsaw,
    colback=white,
    frame hidden, boxrule=0pt, arc=1.8mm,
    colbacktitle=colElec!78!black,
    coltitle=white,
    fonttitle=\sffamily\bfseries\small,
    left=7pt, right=7pt, top=4pt, bottom=4pt,
    toptitle=2.5pt, bottomtitle=2.5pt,
    before skip=8pt, after skip=8pt,
    borderline={0.6pt}{0pt}{colElec!12!white},
  }
}

% ── Style encadré jury ───────────────────────────────────────────────
% Bandeau ambre, fond crème, sans cadre.
\tcbset{
  jury/.style={
    breakable, enhanced jigsaw,
    colback=colJuryBg,
    frame hidden, boxrule=0pt, arc=1.8mm,
    colbacktitle=colJuryFrame,
    coltitle=white,
    fonttitle=\sffamily\bfseries,
    left=7pt, right=7pt, top=4pt, bottom=4pt,
    toptitle=3pt, bottomtitle=3pt,
    before skip=10pt, after skip=8pt,
  }
}

% ── En-tête / pied de page communs (refonte 2026) ─────────────────────
% #1 code fiche, #2 titre court, #3 couleur domaine
\newcommand{\ficheHeaders}[3]{%
  \fancyhead[L]{\small\sffamily{\bfseries\textcolor{#3}{#1}}\hspace{0.7em}\textcolor{colSlate}{#2}}%
  \fancyhead[R]{\small\sffamily\textcolor{colSlate}{CAPET SII IE\, --- \,Session 2026}}%
  \fancyfoot[C]{\small\sffamily\textcolor{colSlate}{\thepage\,/\,\pageref{LastPage}}}%
  \fancyhead[C]{}\fancyfoot[L]{}\fancyfoot[R]{}%
}

% ── Bloc en-tête de fiche (BLOC 0, refonte 2026) ──────────────────────
% #1 code  #2 titre  #3 couleur domaine  #4 domaine  #5 ligne présence  #6 sous-thèmes
\newcommand{\ficheEntete}[6]{%
  \begin{tcolorbox}[enhanced, breakable,
    frame hidden, boxrule=0pt, arc=2.5mm,
    interior style={left color=#3, right color=#3!75!black},
    left=9pt, right=9pt, top=7pt, bottom=7pt,
    before skip=0pt, after skip=6pt]
    {\color{white}\sffamily\bfseries\LARGE #1\quad #2}\\[4pt]
    {\color{#3!22!white}\sffamily\small \MakeUppercase{#4}\quad\textbar\quad #5}\\[1pt]
    {\color{#3!35!white}\sffamily\footnotesize Sous-thèmes : #6}
  \end{tcolorbox}%
}

% ── Mise en page ─────────────────────────────────────────────────────
\pagestyle{fancy}
\fancyhf{}
\renewcommand{\headrulewidth}{0.5pt}
\renewcommand{\footrulewidth}{0pt}
\AtBeginDocument{%
  \renewcommand{\headrule}{{\color{colRule}\hrule width\headwidth height 0.5pt}}%
}

\setlist[itemize]{leftmargin=1.2em, itemsep=1pt, topsep=2pt, parsep=0pt}
\setlist[enumerate]{leftmargin=1.2em, itemsep=3pt, topsep=2pt, parsep=0pt}

\setlength{\parskip}{2pt}
\setlength{\columnsep}{10pt}

\def\CAPETBEGINFILE{\begin{document}}
\def\CAPETENDFILE{\end{document}}
\else
\def\CAPETBEGINFILE{}
\def\CAPETENDFILE{}
\fi

\ficheHeaders{T15}{Stockage d'énergie}{colElec}

\CAPETBEGINFILE

% -----------------------------------------------------------------------
%  BLOC EN-TETE
% -----------------------------------------------------------------------

\ficheEntete{T15}{Stockage d'énergie électrique}{colElec}%
  {électrotechnique - Systèmes énergétiques}%
  {\textbf{Présence :} \textbf{5/13} sujets (2013-2026) - thème \emph{régulier}}%
  {capacité batterie (A\,h) - association série/parallèle de cellules - taux de décharge (DoD) - autonomie - STEP micro-hydraulique - technologies Li-Ion / NMC / LFP - BMS et SOC}

% -----------------------------------------------------------------------
%  1. RAPPELS DE COURS
% -----------------------------------------------------------------------

\begin{tcolorbox}[
  ficheSection=colElec,
  title={1. Rappels de cours},
]

% -- 1.1 Capacite et energie d'une batterie ---------------------------
\textbf{1.1 Capacité et énergie d'une batterie}

\smallskip
\[
  E = U_{nom} \times C_b
  \qquad;\qquad
  C_b = \frac{E_{necessaire}}{U_{nom}}
  \qquad;\qquad
  E_{utile} = E_{totale} \times (1 - DoD_{min})
\]
$C_b$ en A\,h, $E$ en W\,h, $U_{nom}$ en V. Le taux de décharge minimum $DoD_{min}$ protège les cellules (ex : $DoD_{min} = 20\,\%$ - énergie utile = 80\,\% de l'énergie totale).

\smallskip
\textbf{Attention aux deux conventions rencontrées dans les sujets} : « ne pas excéder 70\,\% de décharge » (rapport 2020) $\Rightarrow$ énergie utile $= 70\,\%$ de $E_{totale}$ ($E_{necessaire} = E_{consommee}/0{,}7$) ; « exploitée jusqu'à 20\,\% de décharge » (niveau critique, sujet 2026) $\Rightarrow$ énergie utile $= 80\,\%$. Lire précisément l'énoncé -- le corrigé 2020 avertit : «~Attention à ne pas confondre \% de décharge et profondeur de décharge\ldots~»

\medskip\hrule\vspace{4pt}

% -- 1.2 Association de cellules ---------------------------------------
\textbf{1.2 Association série / parallèle de cellules}

\smallskip
\begin{center}
\renewcommand{\arraystretch}{1.4}
\begin{tabular}{llll}
\toprule
\textbf{Association} & \textbf{Tension résultante} & \textbf{Capacité résultante} & \textbf{Formule} \\
\midrule
Série & $N_s \times U_{cell}$ & $C_{cell}$ inchangée & $N_s = U_{bat}/U_{cell}$ \\
Parallèle & $U_{cell}$ inchangée & $N_p \times C_{cell}$ & $N_p = C_{bat}/C_{cell}$ \\
\bottomrule
\end{tabular}
\end{center}
$N_{total} = N_s \times N_p$.

\smallskip
\begin{center}
\begin{tikzpicture}[scale=0.9, font=\footnotesize]
  \draw (0,0) rectangle (0.9,0.5); \draw (0,0.6) rectangle (0.9,1.1); \draw (0,1.2) rectangle (0.9,1.7);
  \draw (1.2,0) rectangle (2.1,0.5); \draw (1.2,0.6) rectangle (2.1,1.1); \draw (1.2,1.2) rectangle (2.1,1.7);
  \draw (2.4,0) rectangle (3.3,0.5); \draw (2.4,0.6) rectangle (3.3,1.1); \draw (2.4,1.2) rectangle (3.3,1.7);
  \node[left] at (-0.15,0.85) {$N_s$};
  \node[below] at (1.65,-0.2) {$N_p$ branches};
\end{tikzpicture}
\end{center}
Chaque colonne empile $N_s$ cellules en série (tensions cumulées) ; les $N_p$ colonnes sont mises en parallèle (capacités cumulées).

\medskip\hrule\vspace{4pt}

% -- 1.3 Autonomie -----------------------------------------------------
\textbf{1.3 Autonomie et calcul de capacité}

\smallskip
\[
  t_{autonomie} = \frac{E_{utile}}{P_{consommee}}
  \quad\text{(h)}
  \qquad;\qquad
  d_{autonomie} = t_{autonomie} \times v
  \quad\text{(km)}
\]

\medskip\hrule\vspace{4pt}

% -- 1.4 STEP - stockage hydraulique ----------------------------------
\textbf{1.4 STEP micro-hydraulique (Station de Transfert d'Énergie par Pompage)}

\smallskip
\[
  E_p = m\,g\,h = \rho_{eau}\,V_{utile}\,g\,h
  \qquad;\qquad
  E_{elec} = E_p \times \eta_{turbinage}
  \qquad;\qquad
  \eta_{st} = \frac{E_{elec\,turbinage}}{E_{elec\,pompage}}
\]
$\rho_{eau} = 1000\,\mathrm{kg\!\cdot\!m^{-3}}$, $g = 9{,}81\,\mathrm{m\!\cdot\!s^{-2}}$.
Durée de remplissage : $T_{remplissage} = V_{utile}/Q_{pompe}$.

\medskip\hrule\vspace{4pt}

% -- 1.5 Technologies de cellules -------------------------------------
\textbf{1.5 Technologies de cellules}

\smallskip
\begin{center}
\renewcommand{\arraystretch}{1.2}
\begin{tabular}{llll}
\toprule
\textbf{Technologie} & \textbf{$U_{nom}$} & \textbf{Énergie massique} & \textbf{Application typique} \\
\midrule
Li-Ion / NMC & 3{,}5-3{,}7\,V & 150-250 W\,h\,kg$^{-1}$ & Bus, VE, robot \\
Li-FePO4 (LFP) & 3{,}2\,V & 100-160 W\,h\,kg$^{-1}$ & Stockage stationnaire \\
Plomb-acide & 2{,}0\,V & 30-50 W\,h\,kg$^{-1}$ & Démarrage, bateau \\
Supercondensateur & 2{,}5-3{,}0\,V & 5-20 W\,h\,kg$^{-1}$ & Tampon puissance \\
\bottomrule
\end{tabular}
\end{center}

\end{tcolorbox}

% -----------------------------------------------------------------------
%  2. FORMULAIRE SYNTHETIQUE
% -----------------------------------------------------------------------

\begin{tcolorbox}[
  ficheSection=colElec,
  title={2. Formulaire synthétique},
]
\small
\begin{multicols}{2}
\renewcommand{\arraystretch}{1.45}
\begin{tabular}{@{}p{4.5cm}ll@{}}
\toprule
\textbf{Formule} & \textbf{Usage} & \textbf{U.} \\
\midrule
\multicolumn{3}{@{}l}{\footnotesize\textit{Batterie}} \\
$E = U_{nom} \cdot C_b$ & Énergie & W\,h \\
$C_b = E/U_{nom}$ & Capacité & A\,h \\
$E_{utile} = E(1-DoD_{min})$ & Utile & W\,h \\
$t = E_{utile}/P$ & Autonomie & h \\
$N_s = U_{bat}/U_{cell}$ & Série & - \\
$N_p = C_{bat}/C_{cell}$ & Parallèle & - \\
\bottomrule
\end{tabular}

\columnbreak

\renewcommand{\arraystretch}{1.45}
\begin{tabular}{@{}p{4.5cm}ll@{}}
\toprule
\textbf{Formule} & \textbf{Usage} & \textbf{U.} \\
\midrule
\multicolumn{3}{@{}l}{\footnotesize\textit{STEP}} \\
$E_p = \rho V g h$ & Énergie pot. & J \\
$E_{elec} = E_p \cdot \eta_{turb}$ & Turbinage & W\,h \\
$\eta_{st} = E_{turb}/E_{pomp}$ & Rendement & - \\
$T = V_{utile}/Q$ & Remplissage & h \\
\midrule
\multicolumn{3}{@{}l}{\footnotesize\textit{Données typiques}} \\
$U_{NMC} = 3{,}5\,\mathrm{V}$ & Urbanloop & - \\
$\eta_{STEP} \approx 58\,\%$ & Micro-STEP & - \\
\bottomrule
\end{tabular}
\end{multicols}
\normalsize
\end{tcolorbox}

% -----------------------------------------------------------------------
%  3. QUESTIONS TYPES
% -----------------------------------------------------------------------

\begin{tcolorbox}[
  ficheSection=colElec,
  title={3. Questions types - extraites des sujets réels},
]

% -- QT1 -------------------------------------------------------------

\begin{tcolorbox}[questiontype,
  title={QT1 - Association série/parallèle et autonomie - bus Hycity \hfill\normalfont\footnotesize E2 appl. 2026 (Q A3-A5)}
]
\textbf{Énoncé.}
Batterie bus Hycity : $U_{bat} = 600\,\mathrm{V}$, $E_{bat} = 128{,}4\,\mathrm{kW\!\cdot\!h}$, cellules NMC HpCO ($U_{cell} = 3{,}5\,\mathrm{V}$, $C_{cell} = 53{,}5\,\mathrm{A\!\cdot\!h}$), décharge min 20\,\%. Calculer $N_s$ (cellules série par module), $C_{totale}$, puis l'autonomie à $v = 40\,\mathrm{km\!\cdot\!h^{-1}}$ pour $P_{moy} = 22\,\mathrm{kW}$.

\tcbline

\textbf{Corrigé.}

\textbf{Étape 1 — Nombre de cellules en série.}
\[
  N_s = \frac{600/2}{3{,}5} = \frac{300}{3{,}5} \approx 86\,\text{cellules/module} \quad (2\,\text{modules en série} \times 2\,\text{branches})
\]

\textbf{Étape 2 — Capacité totale.}
\[
  C_{totale} = \frac{E_{bat}}{U_{bat}} = \frac{128{,}4 \times 10^3\,\mathrm{W\!\cdot\!h}}{600\,\mathrm{V}} = 214\,\mathrm{A\!\cdot\!h}
\]

\textbf{Étape 3 — Autonomie.}
\[
  E_{utile} = 128{,}4 \times (1 - 0{,}20) = 102{,}72\,\mathrm{kW\!\cdot\!h}
  \quad;\quad
  t = \frac{102{,}72}{22} = 4{,}67\,\mathrm{h}
  \quad;\quad
  d = 4{,}67 \times 40 = \mathbf{186{,}8\,\mathrm{km}}
\]

\textbf{Conclusion.} Avec 186,8\,km d'autonomie, le bus peut rentrer au dépôt depuis n'importe quel point du trajet type.\checkmark
\end{tcolorbox}

\vspace{4pt}

% -- QT2 -------------------------------------------------------------

\begin{tcolorbox}[questiontype,
  title={QT2 - Dimensionnement batterie - taux de décharge et nombre de cellules \hfill\normalfont\footnotesize Rapport jury 2020 (Q20-Q22)}
]
\textbf{Énoncé.}
Un robot effectue 160 trajets types (énergie consommée par trajet type $E_{trajet} \approx \SI{4.06}{\watt\hour}$ -- ne pas confondre avec l'énergie récupérée en descente, $\SI{0.6}{\watt\hour}$, Q16). Taux de décharge max : 70\,\%. Tension nominale batterie : \SI{48}{\volt}. Cellules Li-Ion : $U_{cell} = \SI{3.7}{\volt}$, $C_{cell} = \SI{5.2}{\ampere\hour}$, masse \SI{92}{\gram}. Calculer $E_{necessaire}$ puis $C_b$, la structure série/parallèle, et vérifier la masse (limite \SI{5}{\kilo\gram}).

\tcbline

\textbf{Corrigé.}

\textbf{Étape 1 — Énergie nécessaire et capacité.}
\[
  E_{necessaire} = \frac{160 \times 4{,}06}{0{,}7} \approx 928\,\si{\watt\hour}
  \qquad;\qquad
  C_b = \frac{928}{48} = 19{,}33\,\si{\ampere\hour}
\]
$C_b = \SI{19.33}{\ampere\hour}$ est la valeur du corrigé officiel (rapport jury 2020, Q21-Q22) ; $E_{trajet} \approx \SI{4.06}{\watt\hour}$ est [corrigé reconstitué] par calcul inverse, la valeur exacte de Q19 n'étant pas lisible dans la source.

\textbf{Étape 2 — Structure série/parallèle.}
\[
  N_s = \frac{48}{3{,}7} = 12{,}97 \to 13\,\text{cellules en série}
  \quad;\quad
  N_p = \frac{19{,}33}{5{,}2} = 3{,}72 \to 4\,\text{branches en parallèle}
\]

\textbf{Étape 3 — Vérification masse.}
$N_{total} = 13 \times 4 = 52\,\text{cellules}$. Masse = $52 \times 0{,}092 = 4{,}784\,\si{\kilo\gram}$.

\textbf{Conclusion.} $\SI{4.784}{\kilo\gram} < \SI{5}{\kilo\gram}$ : l'exigence « masse de la batterie $< \SI{5}{\kilo\gram}$ » du cahier des charges est respectée, avec une marge d'environ 4\,\% qu'il convient de signaler.\checkmark
\end{tcolorbox}

\vspace{4pt}

% -- QT3 -------------------------------------------------------------

\begin{tcolorbox}[questiontype,
  title={QT3 - Énergie potentielle et rendement STEP \hfill\normalfont\footnotesize E2 disc. 2021 (Q5-Q10), Rapport jury 2021}
]
\textbf{Énoncé.}
Réservoir : volume total 1000\,m$^3$, volume utile 90\,\%. Hauteur de chute $h = 130\,\mathrm{m}$. Rendement turbinage : 78\,\%. Puissance pompe : 76\,kW, durée remplissage : 5,625\,h. Calculer $E_p$, $E_{elec\,turbinage}$ et le rendement global $\eta_{st}$.

\tcbline

\textbf{Corrigé.}

\textbf{Étape 1 — Énergie potentielle.}
\[
  V_{utile} = 0{,}90 \times 1000 = 900\,\mathrm{m^3}
  \quad;\quad
  m = 1000 \times 900 = 9 \times 10^5\,\mathrm{kg}
\]
\[
  E_p = m\,g\,h = 9 \times 10^5 \times 9{,}81 \times 130 = 1{,}147 \times 10^9\,\mathrm{J} = \frac{1{,}147 \times 10^9}{3{,}6 \times 10^6} = \mathbf{318{,}8\,\mathrm{kW\!\cdot\!h}}
\]

\textbf{Étape 2 — Énergie électrique turbinée.}
\[
  E_{elec\,turbinage} = 318{,}8 \times 0{,}78 = \mathbf{248{,}7\,\mathrm{kW\!\cdot\!h}}
\]

\textbf{Étape 3 — Rendement du cycle complet.}
\[
  E_{pompage} = 76 \times 5{,}625 = 427{,}5\,\mathrm{kW\!\cdot\!h}
  \quad;\quad
  \eta_{st} = \frac{248{,}7}{427{,}5} = \mathbf{58{,}2\,\%}
\]

\textbf{Conclusion.} Un rendement de 58,2\,\% est faible comparé à une batterie électrochimique, mais la technologie reste écologique et complémentaire aux panneaux PV (production stockée pendant les heures creuses pour restitution la nuit).
\end{tcolorbox}

\end{tcolorbox}

% -----------------------------------------------------------------------
%  4. ATTENTES DU JURY
% -----------------------------------------------------------------------

\begin{tcolorbox}[
  jury,
  title={4. Attentes du jury - Erreurs récurrentes et points de vigilance},
]

\begin{enumerate}[label=\textbf{\arabic*.}, leftmargin=*, itemsep=5pt]

  \item \textbf{Confusion W\,h et A\,h.}
  La capacité s'exprime en A\,h, l'énergie en W\,h. La conversion est $E = U_{nom} \times C_b$. Ne pas additionner des A\,h et des W\,h.

  \item \textbf{Oubli du taux de décharge.}
  Si le constructeur impose $DoD_{max} = 70\,\%$, l'énergie utile n'est que 70\,\% de l'énergie totale. Rapport jury 2020 : «~Une critique à la fois qualitative et quantitative était attendue afin de vérifier que l'installation permette de ne pas excéder 70~\% de décharge de la batterie.~»

  \item \textbf{Erreur sur l'association série/parallèle.}
  En \emph{série} : les tensions s'additionnent, la capacité reste celle d'une cellule. En \emph{parallèle} : les capacités s'additionnent, la tension reste celle d'une cellule. Ne jamais multiplier les deux simultanément.

  \item \textbf{Densité de l'eau dans la formule STEP.}
  Utiliser $E_p = \rho V g h$ avec $\rho_{eau} = 1000\,\mathrm{kg\!\cdot\!m^{-3}}$. Oublier la densité donne un résultat 1000 fois trop faible.

  \item \textbf{Rendement STEP = cycle complet.}
  Calculer $\eta_{st} = E_{turbinage}/E_{pompage}$ sur le cycle complet, pas uniquement sur la phase de turbinage (qui donnerait 100\,\%). Citer la valeur numérique ($\approx$58\,\%) ET la conclusion contextuelle (écologique, peu performant).

  \item \textbf{Réponses qualitatives sans chiffres.}
  Rapport jury 2020 : «~Le jury déplore des réponses trop évasives sans aucun lien avec les documents techniques et/ou un manque de valeurs numériques pour assoir les justifications.~» Toujours appuyer la conclusion (autonomie suffisante) par un calcul comparant $E_{disponible}$ et $E_{necessaire}$.

\end{enumerate}

\medskip
\begin{tcolorbox}[
  colback=colJuryFrame!12!colJury,
  colframe=colJuryFrame,
  left=4pt, right=4pt, top=2pt, bottom=2pt,
  fontupper=\small\itshape
]
Rapport jury 2020 : «~Le jury note ici les progrès des candidats quant aux conversions d'unités mais déplore que des problématiques simples, comme le dimensionnement d'une batterie, ne soient pas maîtrisées.~»\\
Rapport jury 2021 (corrigé officiel, Q10) : «~Le procédé de stockage par pompage n'est pas très performant (rendement de 58\,\%). [\ldots] Mais ce procédé est plus respectueux de l'environnement que le stockage par batterie.~»
\end{tcolorbox}

\end{tcolorbox}

% -----------------------------------------------------------------------
%  PIED - PRESENCE DANS LES SUJETS
% -----------------------------------------------------------------------

\begin{tcolorbox}[
  ficheSection=colPresence,
  title={Présence dans les sujets analysés (2013-2026)},
]
\begin{center}
\footnotesize\renewcommand{\arraystretch}{1.15}
\begin{tabular}{lccccccccccccc}
\toprule
\textbf{Année} & 13 & 14 & 15 & 16 & 17 & 18 & 19 & 20 & 21 & 22 & 23 & 24 & 25 \\
\midrule
E1 (commune)   & - & - & - & - & - & - & - & - & - & - & - & - & - \\
E2 discipl. IE & - & - & - & - & - & - & - & $\checkmark$ & $\checkmark$ & - & $\checkmark$ & $\checkmark$ & $\checkmark$ \\
E2 appl. comm. & - & - & - & - & - & - & - & - & - & - & - & - & - \\
\bottomrule
\end{tabular}
\end{center}
\smallskip
\footnotesize\textbf{Note :} la session 2026 inclut un sujet BAC3 (bus électrique Hycity) qui traite le stockage Li-Ion 600\,V / 128{,}4\,kW\,h en détail (Q A3-A5).\\
\textbf{Sous-thèmes fréquents :}
capacité batterie en A\,h (5/13) - association série/parallèle (4/13) - taux de décharge DoD (4/13) - énergie potentielle STEP (2/13) - rendement cycle complet (2/13)

\medskip
\textbf{Ordres de grandeur (repères rapides).}
\begin{itemize}[leftmargin=1.3em]
  \item Réseau BT FR : tension phase-neutre typique \SI{230}{\volt}, tension phase-phase typique \SI{400}{\volt}, fréquence \SI{50}{\hertz}.
  \item Courants d'étude fréquents en sujet : de quelques \si{\ampere} à quelques dizaines de \si{\ampere}.
  \item En conversion d'énergie, un rendement attendu reste généralement entre 0{,}8 et 0{,}98 selon la technologie.
\end{itemize}

\medskip
\textbf{Synthèse express (à refaire sans support).}
\begin{enumerate}[leftmargin=1.4em]
  \item Poser les hypothèses du modèle et les conventions de signe.
  \item Écrire 2 à 4 relations clés avec unités.
  \item Réaliser un calcul numérique complet et vérifier l'ordre de grandeur.
  \item Justifier le résultat par un argument technique de niveau CAPET.
\end{enumerate}

\end{tcolorbox}

\CAPETENDFILE
